\chapter{Criterio de arquitectura externa final del corpus}

\section{Arquitectura externa final}
\begin{definition}
Se llama arquitectura externa final del corpus al orden disciplinado mediante el cual MADELEINE puede presentarse públicamente como sistema, conservando la diferencia funcional entre fundamento, canon intermedio, capa metodológica, claims públicos, material observado y reservas internas.
\end{definition}

\begin{observation}
La arquitectura externa final no equivale a mera suma de materiales visibles. Su función es fijar una forma pública coherente del sistema sin borrar su jerarquía interna ni deformar su proceso metodológico.
\end{observation}

\section{Condición mínima de arquitectura}
\begin{definition}
Existe arquitectura externa final del corpus cuando pueden declararse simultáneamente:
\begin{enumerate}
  \item una relación legible entre Tomo I, Tomo II y Tomo III;
  \item una capa pública sobria compuesta por fichas, dossiers, mapas o salidas integrales compatibles;
  \item una distinción suficiente entre material publicable, material observado y material todavía reservado;
  \item trazabilidad visible entre lo expuesto externamente y su genealogía metodológica interna;
  \item ausencia de colapso entre fundamento, organización operativa y proyección pública.
\end{enumerate}
\end{definition}

\begin{observation}
No se exige todavía una clausura definitiva de toda evolución futura del sistema. Basta con que la forma externa ya no contradiga la estructura interna alcanzada.
\end{observation}

\section{Arquitectura y tomos}
\begin{definition}
Toda arquitectura externa final debe preservar la función propia de cada tomo:
\begin{enumerate}
  \item Tomo I como fijación del núcleo y del canon basal;
  \item Tomo II como canon intermedio y régimen de organización operativa;
  \item Tomo III como capa metodológica, curatorial y de proyección disciplinada.
\end{enumerate}
\end{definition}

\begin{observation}
Esta cláusula impide que la presentación externa convierta los tres tomos en un relato indiferenciado o en una colección sin jerarquía.
\end{observation}

\section{Arquitectura y capa pública}
\begin{definition}
La arquitectura externa final puede articular:
\begin{enumerate}
  \item fichas públicas mínimas;
  \item dossiers externos de claims seleccionados;
  \item mapa público mínimo del corpus;
  \item publicación integral sobria del sistema;
  \item referencias explícitas a material observado o reservado cuando ello resulte necesario para mantener honestidad metodológica.
\end{enumerate}
\end{definition}

\begin{observation}
La capa pública no debe simular completitud absoluta allí donde el propio sistema preserva reservas, gradientes de estabilidad o material aún no proyectable.
\end{observation}

\section{Condición mínima de no arquitectura}
\begin{definition}
No debe declararse arquitectura externa final cuando la presentación pública dependa de simplificaciones que destruyan la diferencia entre núcleo, canon, metodología, capa pública y reserva interna.
\end{definition}

\begin{observation}
La no declaración de arquitectura final no implica insuficiencia del sistema. Puede indicar solo que aún conviene mantener mediaciones parciales antes de cerrar una forma externa total.
\end{observation}

\section{Cláusula de alcance}
\begin{observation}
El criterio de arquitectura externa final del corpus fija la primera condición de forma pública completa y disciplinada del sistema MADELEINE. No constituye todavía una política editorial definitiva ni una clausura total de toda expansión futura del corpus.
\end{observation}

